%======================================================================
\section{The model}\label{sec:model}
%======================================================================

\subsection{Probability space and measures}

Fix a filtered probability space $(\Omega,\F,(\F_t)_{t\ge0},\Q)$ satisfying the usual
conditions, supporting two standard Brownian motions $W$ and $B$ with
$\dd\langle W,B\rangle_t=\rho\dd t$, a Poisson process $N$ with intensity $\lambda_J$,
an i.i.d.\ sequence of jump sizes $(J_n)_{n\ge1}$, and two unit exponential random
variables $E_{\mathrm{res}}$ and $E_{\mathrm{spec}}$. The Brownian motions, the
compound Poisson process and the exponentials are mutually independent, except for the
correlation between $W$ and $B$. Denote by $\G=(\G_t)$ the filtration generated by
$(W,B,N,(J_n))$, the information in the state paths, and by
$\G_\infty=\bigvee_t\G_t$ its terminal $\sigma$-algebra.

All prices are expectations under $\Q$, a pricing measure. The CET1 ratio is not a
traded asset and its dynamics are estimated from quarterly history, that is under the
physical measure $\Prob$.

\begin{assumption}[Capital dynamics under the pricing measure]\label{ass:measure}
The dynamics of the CET1 ratio are the same under $\Q$ and under $\Prob$. All risk
premia not captured by the mechanical features of the contract are absorbed by the
write-down hazard $\lambda$, which is therefore a risk-neutral intensity and not a
probability of default.
\end{assumption}

Assumption~\ref{ass:measure} is the standard compromise of reduced-form models
calibrated to an untraded state variable. Its consequence is that $\lambda$ should be
read as a price, not as a forecast: a hazard of $2.2\%$ a year is the compensation the
market requires for the write-down risk and its premium, not a $2.2\%$ annual
probability of write-down.

\subsection{The CET1 ratio}

Let $C_t$ be the CET1 ratio in percentage points. It follows the mean-reverting
jump-diffusion
\begin{equation}\label{eq:cet1}
  \dd C_t = \kappa(\theta-C_t)\dd t + \sigma\dd W_t - \dd L_t,
  \qquad L_t=\sum_{n=1}^{N_t}J_n,
\end{equation}
with $\kappa>0$, $\sigma>0$, and lognormal jump sizes with mean $m_J$,
\begin{equation}\label{eq:jumps}
  J_n = m_J\exp\!\bigl(\sigma_J Z_n-\tfrac12\sigma_J^2\bigr),\qquad Z_n\sim N(0,1),
  \qquad \E[J_n]=m_J,\quad \E[J_n^2]=m_J^2e^{\sigma_J^2}.
\end{equation}
The drift represents capital management: retained earnings, distribution policy and
the management of risk-weighted assets pull the ratio back to the target $\theta$, the
CET1 objective the bank communicates to investors. Jumps are one-sided because losses
(a fine, a trading loss, an impairment, a regulatory change in risk weights) arrive
suddenly, while capital is rebuilt gradually.

\begin{proposition}[Solution and moments]\label{prop:cet1}
Equation~\eqref{eq:cet1} has the unique strong solution
\begin{equation}\label{eq:cet1sol}
  C_t = \theta + (C_0-\theta)e^{-\kappa t}
  + \sigma\int_0^t e^{-\kappa(t-u)}\dd W_u
  - \int_0^t e^{-\kappa(t-u)}\dd L_u ,
\end{equation}
with mean and variance
\begin{align}
  \E[C_t] &= \theta + (C_0-\theta)e^{-\kappa t} - \frac{\lambda_J m_J}{\kappa}\bigl(1-e^{-\kappa t}\bigr),
  \label{eq:cet1mean}\\
  \operatorname{Var}[C_t] &= \frac{\sigma^2+\lambda_J m_J^2e^{\sigma_J^2}}{2\kappa}\bigl(1-e^{-2\kappa t}\bigr).
  \label{eq:cet1var}
\end{align}
As $t\to\infty$, $C_t$ converges in distribution to a stationary law with mean
$\theta-\lambda_Jm_J/\kappa$ and variance $(\sigma^2+\lambda_Jm_J^2e^{\sigma_J^2})/(2\kappa)$.
\end{proposition}

\begin{proof}
The coefficients are affine with Lipschitz drift and constant diffusion, so a unique
strong solution exists. Apply It\^o's formula for semimartingales to
$e^{\kappa t}C_t$: $\dd(e^{\kappa t}C_t)=e^{\kappa t}(\kappa\theta\dd t+\sigma\dd W_t-\dd L_t)$,
and integrate to obtain \eqref{eq:cet1sol}. For the mean, the stochastic integral against
$W$ has zero expectation, and
$\E\int_0^te^{-\kappa(t-u)}\dd L_u=\lambda_Jm_J\int_0^te^{-\kappa(t-u)}\dd u$ by the
compensator of the compound Poisson process. For the variance, the Brownian and jump
integrals are independent; It\^o's isometry gives
$\sigma^2\int_0^te^{-2\kappa(t-u)}\dd u$ for the first, and the isometry for compensated
compound Poisson integrals gives $\lambda_J\E[J^2]\int_0^te^{-2\kappa(t-u)}\dd u$ for the
second. Evaluating the integrals gives \eqref{eq:cet1mean}--\eqref{eq:cet1var}. The
characteristic function of \eqref{eq:cet1sol} converges as $t\to\infty$, which gives
the stationary law.
\end{proof}

\begin{remark}[Target and long-run mean]\label{rem:target}
The long-run mean is not $\theta$ but $\theta-\lambda_Jm_J/\kappa$: jumps pull the ratio
down on average, and mean reversion offsets them only partially. With the base-case
parameters ($\theta=13\%$, $\kappa=0.35$, $\lambda_J=0.12$, $m_J=2.2$pp) the stationary
mean is $12.25\%$, and the stationary standard deviation is $1.17$pp, of which the
diffusion alone contributes $\sigma/\sqrt{2\kappa}=0.56$pp. Interpreting $\theta$ as
the bank's target is therefore accurate in the absence of shocks; in the presence of
shocks, the model's capital hovers about $0.75$pp below target on average. A calibration
that wanted the long-run mean, rather than the shock-free target, to equal $13\%$ would set
$\theta=13.75\%$; we keep $\theta$ at the communicated target and report the implication.
\end{remark}

The mechanical trigger time is the first passage
\begin{equation}
  \tau_m = \inf\{t\ge0:\ C_t\le \ctrig\},\qquad \ctrig=5.125\%.
\end{equation}

\subsection{The refinancing spread}

Let $S_t$ be the market component of the issuer's new-issue AT1 spread, in decimal.
It follows the Ornstein--Uhlenbeck process
\begin{equation}\label{eq:spread}
  \dd S_t = a(\bar s-S_t)\dd t + \eta\dd B_t,
\end{equation}
correlated with the diffusive capital shocks through $\rho<0$: when capital falls,
spreads tend to widen. The issuer's effective spread adds a capital premium once the
CET1 ratio enters the buffer,
\begin{equation}\label{eq:seff}
  s_t = S_t + \beta\bigl(\cmda(t)-C_t\bigr)^+ ,
\end{equation}
where $\cmda(t)$ is the MDA threshold defined in \eqref{eq:mdathreshold} below and $\beta$
is the extra spread per point of CET1 inside the buffer. The effective spread is the rate
at which the bank could issue a new AT1 at time $t$.

\begin{proposition}[Spread distribution]\label{prop:spread}
The solution of \eqref{eq:spread} is Gaussian,
\begin{equation}
  S_t = \bar s + (S_0-\bar s)e^{-at} + \eta\int_0^te^{-a(t-u)}\dd B_u
  \ \sim\ N\!\Bigl(\bar s + (S_0-\bar s)e^{-at},\ \frac{\eta^2}{2a}\bigl(1-e^{-2at}\bigr)\Bigr),
\end{equation}
and its covariance with the diffusive part of capital,
$C^W_t=\sigma\int_0^te^{-\kappa(t-u)}\dd W_u$, is
\begin{equation}\label{eq:covsc}
  \operatorname{Cov}\bigl(S_t,C^W_t\bigr) = \frac{\rho\sigma\eta}{a+\kappa}\Bigl(1-e^{-(a+\kappa)t}\Bigr).
\end{equation}
In the stationary regime $\Q(S_t<0)=\Phi\bigl(-\bar s\sqrt{2a}/\eta\bigr)$.
\end{proposition}

\begin{proof}
The solution follows as in Proposition~\ref{prop:cet1}. The covariance follows from
It\^o's isometry for two Wiener integrals:
$\E[\int_0^tf\dd B\int_0^tg\dd W]=\rho\int_0^tf(u)g(u)\dd u$ with
$f(u)=\eta e^{-a(t-u)}$ and $g(u)=\sigma e^{-\kappa(t-u)}$.
\end{proof}

With the base-case parameters ($a=0.5$, $\bar s=275$bp, $\eta=100$bp) the stationary
standard deviation of $S$ is $100$bp and the stationary probability of a negative spread
is $0.3\%$. Because the hazard and the call rule use the spread, negative values are
handled explicitly: the hazard is floored at zero (equation~\eqref{eq:hazard}), and a
negative refinancing spread simply triggers a call. A square-root or log-normal spread
would avoid negative values at the cost of the exact Gaussian transition; at a
probability of $0.3\%$ the difference is immaterial.

\subsection{The regulatory threshold}

Let $b$ be the MDA floor (Pillar~1 plus the CET1 part of Pillar~2R) and $w(t)$ the
combined buffer at time $t$. The MDA threshold is
\begin{equation}\label{eq:mdathreshold}
  \cmda(t) = b + w(t),\qquad w(t)=w_0+\sum_k \Delta w_k\,\1\{t\ge t_k\},
\end{equation}
where the steps $(t_k,\Delta w_k)$ are scheduled buffer changes. For BNP, $b=5.64\%$,
$w_0=4.79\%$ (so that $\cmda(0)$ equals the published requirement of $10.43\%$), and the
O-SII step-ups give $\Delta w=0.25$pp on 1 January 2027 and on 1 January 2028.

\subsection{Contract features as functionals of the state}

\paragraph{MDA coupon restriction.} With $q_t=(C_t-b)/w(t)$ the position of the ratio
inside the buffer, the fraction of the scheduled coupon that can be paid is
\begin{equation}\label{eq:mda}
  \phi(C_t,t)=
  \begin{cases}
    1 & C_t\ge \cmda(t),\\
    0.6 & q_t\in[0.75,1),\\
    0.4 & q_t\in[0.50,0.75),\\
    0.2 & q_t\in[0.25,0.50),\\
    0 & q_t<0.25.
  \end{cases}
\end{equation}
This applies the MDA quartile factors to the coupon itself. It is a stylised and
conservative version of the rule, which caps all distributions together
(Section~\ref{sec:institutions}); the resulting MDA component of the spread is an upper
bound.

\paragraph{Call and extension.} Let $T_1<T_2<\dots$ be the reset dates and $m$ the reset
margin. The issuer calls at
\begin{equation}\label{eq:call}
  \tau_c = \min\bigl\{T_k:\ m> s_{T_k}-\delta,\ \ C_{T_k}\ge\cmda(T_k),\ \ T_k<\tau_m\bigr\}
  \qquad(\min\emptyset=+\infty),
\end{equation}
that is, at the first reset date at which keeping the bond costs more than refinancing
it, less a reputational tolerance $\delta\ge0$, provided that the regulator would approve
the call (CET1 above the threshold) and the bond is still alive. In the base case
$\delta=0$. Since each condition is $\G_{T_k}$-measurable, $\tau_c$ is a $\G$-stopping
time taking values in $\{T_k\}\cup\{+\infty\}$. If $\eta>0$ the call is random even for a
healthy bank, and the bond with the lowest margin is the first to be extended when
spreads widen.

\paragraph{Conversion.} At $\tau_m$ the bond converts into shares at the higher of the
market price and the floor price $P_{\mathrm{floor}}$. When the share trades below the
floor, as it will at the trigger, the investor receives
$F\cdot S^{\mathrm{eq}}_{\tau_m}/P_{\mathrm{floor}}$. We do not model the share price
jointly with capital; instead the share price at conversion is a fixed fraction $\xi$ of
today's price $S^{\mathrm{eq}}_0$ converted at today's exchange rate $X_0$, so the recovery
at conversion is
\begin{equation}\label{eq:conv}
  R_{\mathrm{conv}}=\min\!\left(1,\ \frac{\xi\,S^{\mathrm{eq}}_0X_0}{P_{\mathrm{floor}}}\right).
\end{equation}
With $\xi=0.2$, a share price of EUR~98.43 and $X_0=1.15$, $R_{\mathrm{conv}}$ is $37\%$
for the 6.875\% and the 7.45\% and $31\%$ for the 7.20\%, whose floor is higher.

\subsection{The write-down hazard}\label{sec:hazard}

The regulator may impose losses before the trigger. We model the PONV time as a
doubly-stochastic (Cox) time with intensity proportional to the effective spread,
\begin{equation}\label{eq:hazard}
  \lambda_t = \lambda\,\frac{s_t^+}{s_0},\qquad
  I_t=\int_0^t\frac{s_u^+}{s_0}\dd u ,
\end{equation}
so that $\lambda$ is today's hazard and $I_t$ is the integrated driver. The hazard rises
when spreads widen or capital falls inside the buffer. The recovery of the AT1 at PONV is
zero.

Linking the hazard to $s_t$ is what makes extension consistent. The issuer extends when
$s_t>m$; with a state-dependent hazard those are exactly the states in which the bond is
riskier, and the extended cash flows are discounted accordingly. With a constant hazard,
the extended coupon $y^{\mathrm{ref}}+m$ would be discounted at a rate that does not reflect
the state, and in states where $m$ exceeds the risk-adjusted discount spread, not calling
would add value for the investor. Section~\ref{sec:robust} shows the size of this effect.

\subsection{Two clocks: resolution and AT1-specific write-down}\label{sec:twoclocks}

To price the AT1 jointly with the bank's Tier~2, decompose the hazard level into
\begin{equation}\label{eq:split}
  \lambda=\lres+\lspec,\qquad \lres,\lspec\ge0,
\end{equation}
and define two write-down times with the same driver,
\begin{equation}\label{eq:clocks}
  \tau_{\mathrm{res}} = \inf\{t:\ \lres I_t\ge E_{\mathrm{res}}\},\qquad
  \tau_{\mathrm{spec}} = \inf\{t:\ \lspec I_t\ge E_{\mathrm{spec}}\}.
\end{equation}
The resolution time $\tau_{\mathrm{res}}$ imposes losses on AT1 and on Tier~2; the
AT1-specific time $\tau_{\mathrm{spec}}$ imposes losses on AT1 alone. The AT1 write-down
time is $\tau_P=\tau_{\mathrm{res}}\wedge\tau_{\mathrm{spec}}$.

\begin{proposition}[Joint survival]\label{prop:clocks}
Conditionally on $\G_\infty$,
\begin{align}
  \Q(\tau_P>t\mid\G_\infty) &= e^{-\lambda I_t},\label{eq:survat1}\\
  \Q(\tau_{\mathrm{res}}>t\mid\G_\infty) &= e^{-\lres I_t},\label{eq:survt2}\\
  \Q(\tau_{\mathrm{res}}\le t,\ \tau_{\mathrm{res}}\le\tau_{\mathrm{spec}}\mid\G_\infty)
  &= \frac{\lres}{\lambda}\bigl(1-e^{-\lambda I_t}\bigr).\label{eq:firstres}
\end{align}
In particular $\tau_P$ is a Cox time with intensity $\lambda s_t^+/s_0$, and the
probability that an AT1 write-down is a resolution event that also hits Tier~2 is
$\lres/\lambda$, independently of the path.
\end{proposition}

\begin{proof}
$I_t$ is $\G_\infty$-measurable, continuous and non-decreasing, and the exponentials are
independent of $\G_\infty$ and of each other. Hence
$\Q(\tau_{\mathrm{res}}>t\mid\G_\infty)=\Q(E_{\mathrm{res}}>\lres I_t\mid\G_\infty)=e^{-\lres I_t}$,
which is \eqref{eq:survt2}; the same holds for $\tau_{\mathrm{spec}}$ with $\lspec$, and by
conditional independence $\Q(\tau_P>t\mid\G_\infty)=e^{-\lres I_t}e^{-\lspec I_t}$,
which is \eqref{eq:survat1}. For \eqref{eq:firstres}, change time to
$u=I_t$: conditionally on $\G_\infty$, $\tau_{\mathrm{res}}$ and $\tau_{\mathrm{spec}}$
become independent exponential times in $u$ with rates $\lres$ and $\lspec$, and the
probability that the first of two independent exponentials occurs before $u$ and is the
one with rate $\lres$ is $\frac{\lres}{\lambda}(1-e^{-\lambda u})$.
\end{proof}

Proposition~\ref{prop:clocks} is the formal content of the Tier~2 anchor. The AT1 price
depends on $\lres$ and $\lspec$ only through their sum, so the AT1 alone cannot separate
them; the Tier~2 depends on $\lres$ alone and identifies it. The ratio $\lres/\lambda$ is
the share of AT1 write-down risk that the market prices in Tier~2 as well.

Two remarks on interpretation. First, $\lres$ is calibrated to the Tier~2 spread, which
also contains the risk and liquidity premia that investors require on Tier~2 itself.
``Priced in Tier~2'' is therefore the accurate label for $\lres$; ``bank resolution
risk'' is an approximation. Second, the construction assumes that resolution hits AT1
whenever it hits Tier~2, which is the creditor hierarchy, and allows AT1-only write-downs
of the Credit Suisse type through $\tau_{\mathrm{spec}}$. It rules out a Tier~2 loss
without an AT1 loss, which the hierarchy also rules out.
